\section{The anticipatory flood signal}
\label{sec:flood-trend}

Flood is the only hazard whose temporal logic does not rest solely on the
validator. A water level has a trend, and the trend is the alert: knowing that
a street is flooded is useful, but knowing that it is filling, and how fast,
is what buys evacuation time.

\paragraph{Not a novel contribution, and read this before describing it as
one.} \Citet{choi2026pagehinkley} target the same gap on the same
infrastructure — anticipatory urban flood warning from existing CCTV, using
Page--Hinkley change detection with short-term trend analysis producing an
estimated time-to-threshold. That is the same capability, published
independently. VIGÍA therefore implements \emph{both} mechanisms and measures
them against each other on identical inputs, rather than claiming either. What
is ours is the multi-hazard pipeline they sit inside.

\subsection{Two signals}

Segmenting water yields a mask $W_t \subseteq \Omega$ over the frame domain
$\Omega$. Two scalars are derived, deliberately:
\begin{align}
  a_t &= \frac{|W_t|}{|\Omega|}
     &&\text{area fraction — needs no calibration}
     \label{eq:area}\\
  \ell_t &= \frac{1}{K}\max\{\, k : \textstyle\sum_{j=k-4}^{k}
              \mathbb{1}[\, p_j \in W_t \,] > 0 \,\}
     &&\text{level along a reference line}
     \label{eq:level}
\end{align}
where $p_1, \dots, p_K$ are $K = 200$ points sampled from the bottom of a
per-camera reference line upward. \Cref{eq:level} walks up from the bottom and
stops at the first sustained dry stretch rather than taking the highest water
pixel anywhere on the line; a reflection or a puddle further up would
otherwise read as a far higher level than reality. The five-sample tolerance
is 2.5\,\% of the line, which absorbs segmentation speckle without jumping
gaps.

The area signal works on any camera the moment it is connected; the level
signal needs a one-off human calibration but is physically meaningful and
comparable between cameras. Choosing the reference line is deliberately not
automated: a wrong automatic choice would silently corrupt every reading from
that camera.

\subsection{Windowed least squares}

Rate of rise is fitted, not differenced. Frame-to-frame differencing of
a noisy segmentation mask produces a rate estimate dominated by segmentation
jitter; an ordinary least-squares slope over a window is stable enough to
threshold on. Over the readings in a trailing window $\mathcal{W}$,
\begin{equation}
  \hat{\beta} = \frac{\sum_{i \in \mathcal{W}} (t_i - \bar{t})(y_i - \bar{y})}
                     {\sum_{i \in \mathcal{W}} (t_i - \bar{t})^{2}},
  \qquad
  R^{2} = 1 - \frac{\sum_i (y_i - \hat{y}_i)^{2}}
                   {\sum_i (y_i - \bar{y})^{2}}.
  \label{eq:ols}
\end{equation}
A trend is declared only when $|\mathcal{W}| \ge 4$ and $R^{2} \ge 0.5$: a
poor fit usually means the segmentation is unstable rather than that the water
is doing something complicated, so the system declines to call it. Where a
calibrated level exists it is preferred over area for the trend decision,
because area fraction moves with camera perspective and level does not. The
operational output is the extrapolation
\begin{equation}
  t_{\text{threshold}} = \frac{\ell^{*} - \ell_t}{\hat{\beta}},
  \qquad \hat{\beta} > 0,
  \label{eq:tttp}
\end{equation}
which is the number an operator actually wants: not ``the street is flooding''
but ``this underpass is impassable in eleven minutes''.

\begin{recorded}{A bug only real data could find}
With a fixed \SI{300}{\second} window, \cref{eq:ols} never fired at all on a
creek camera sampled every fifteen minutes — it reported \textsc{unknown} for
seven hours while the water rose by 10\,\% of the frame — because the window
could never contain two frames. Every synthetic test passed throughout. A time
window is the wrong abstraction when the sampling rate is a property of the
deployment, so the window is now adaptive: it admits at least four sampling
intervals, measured from the observed median gap. Page--Hinkley, being
sample-based, not time-based, was unaffected.
\end{recorded}

\subsection{Page--Hinkley change detection}
\label{sec:ph}

The Page--Hinkley test~\citep{page1954continuous,hinkley1971inference} is a
classical sequential change-point detector. For observations $x_1, x_2, \dots$
with running mean $\bar{x}_i$, tolerated drift $\delta$ and threshold
$\lambda$:
\begin{align}
  m_t &= \sum_{i=1}^{t} \bigl( x_i - \bar{x}_i - \delta \bigr),
  \label{eq:ph-cum}\\
  M_t &= \min_{1 \le i \le t} m_i,
  \label{eq:ph-min}\\
  \mathrm{PH}_t &= m_t - M_t,
  \label{eq:ph-stat}
\end{align}
with an alarm at the first $t \ge t_{\min}$ for which
$\mathrm{PH}_t > \lambda$. Its character differs from \cref{eq:ols}: it
responds quickly to the \emph{onset} of a change, whereas the fit estimates
the \emph{rate} of an ongoing one more stably. They are complementary rather
than competing, which is why both are available.

\subsection{Calibration matters more than the algorithm}
\label{sec:ph-calibration}

On a signal that is not changing, \cref{eq:ph-cum} is a random walk with step
standard deviation $\sigma$ and drift $-\delta$ per sample, and
\cref{eq:ph-stat} is that walk's drawdown below its own running minimum. A
threshold small relative to the noise will therefore be crossed eventually by
noise alone. For Brownian motion with negative drift the all-time maximum
drawdown is exponentially distributed, giving
\begin{equation}
  \Pr[\text{false alarm}] \;\approx\;
  \exp\!\left( -\,\frac{2\,\delta\,\lambda}{\sigma^{2}} \right),
  \label{eq:ph-far}
\end{equation}
and, inverting, a design rule for a camera nobody has calibrated:
\begin{equation}
  \lambda \;\gtrsim\; \frac{\sigma^{2}}{2\delta}\,\ln\frac{1}{\alpha}
  \label{eq:ph-design}
\end{equation}
for a target false-alarm probability $\alpha$.

Our first defaults, $\delta = 0.002$ and $\lambda = 0.02$, looked reasonable
and produced a 92\,\% false-alarm rate on \emph{stable} water with realistic
segmentation noise. \Cref{tab:ph} places the measured grid beside
\cref{eq:ph-far}.

\begin{table}[htbp]
\centering\small
\caption{Page--Hinkley false-alarm rate on stable water, measured over 200
trials of 40 samples, against the drawdown model of \cref{eq:ph-far}. All
values are percentages. The model consistently \emph{under}-predicts — by
roughly a factor of two in the interesting region — because it assumes the
infinite-time limit and treats the running-mean subtraction as exact. It
nevertheless reproduces the ordering across every row and column and locates
the usable/unusable boundary correctly, so \cref{eq:ph-design} is a sound
floor for a new camera and not a substitute for measuring it. Generated by
\texttt{docs/build\_paper.py} from
\texttt{eval/results/page\_hinkley\_calibration.json}.}
\label{tab:ph}
\input{generated/page_hinkley}
\end{table}

The chosen default, $\delta = 0.005$ and $\lambda = 0.10$, trades a little
latency for zero false alarms at typical segmentation noise: detection latency
for a genuine rise is 11, 7 and 4 samples at 0.5\,\%, 1\,\% and 3\,\% per
frame respectively. A camera with noisier segmentation needs a higher
threshold, and \cref{eq:ph-design} says roughly how much higher — but it
should be measured per camera, not assumed.

\subsection{Compared on real cameras}

Both mechanisms were run on identical inputs from LSU creek camera sequences
in the V-FloodNet dataset~\citep{liang2022vfloodnet}: real fixed cameras, real
water, real segmentation noise. The two win on different sampling regimes —
the windowed fit is faster on densely sampled cameras, Page--Hinkley on
sparsely sampled ones, because it counts samples rather than seconds. Neither
is claimed as ours, and the comparison is reported; no winner is
declared.

\subsection{Ground truth the other footage does not have}
\label{sec:gage}

USGS cameras~\citep{usgshivis} sit on instrumented gauges, so every frame
arrives with a measured water level beside it — which no other footage in this
project provides. Over the 50 frames of a real flood event (Walnut Creek at
Raleigh, 8 August 2024, gauge rising to \SI{9.00}{\foot} and receding to
\SI{4.25}{\foot}), segmented water coverage correlates with measured gauge
height at $r = +0.50$.

That is a real but moderate relationship, and it is reported because it
\emph{qualifies} the trend claim instead of supporting it unconditionally.
Coverage tracks the water level between a flood day and a receded day, but
within a day at near-constant level the coverage still varies by more than the
level does. Restricting to midday frames does not improve it ($r = +0.46$ to
$+0.50$ across daylight windows), so the residual is not simply low sun.
Camera-derived coverage is evidence of a trend, not a substitute for a gauge.
This is one event on one camera, and \cref{sec:outstanding} lists repeating it
as outstanding work.
